\documentclass[a4paper,9pt]{extarticle}
\usepackage{geometry}
\usepackage{fancyhdr}
\usepackage{multicol}
\usepackage{enumitem}
\usepackage{titlesec}
\usepackage{xcolor}
\usepackage{hyperref}
\usepackage{tcolorbox}
\usepackage{listings}

% Set page margins
\geometry{margin=0.8in}

% Header and footer
\pagestyle{fancy}
\fancyhf{}
\lhead{Lex Cheat Sheet}
\rfoot{\thepage}

% Set font size
\renewcommand{\baselinestretch}{0.9}

% Title format
\titleformat{\section}{\large\bfseries}{\thesection}{1em}{}
\titleformat{\subsection}{\normalsize\bfseries}{\thesubsection}{1em}{}

% Color for section titles
\definecolor{sectioncolor}{RGB}{0, 51, 102}
\titleformat{\section}[block]{\color{sectioncolor}\large\bfseries}{\thesection}{1em}{}

% Listings
\definecolor{codegreen}{rgb}{0,0.6,0}
\definecolor{codegray}{rgb}{0.5,0.5,0.5}
\definecolor{codepurple}{rgb}{0.58,0,0.82}
\definecolor{backcolour}{rgb}{0.95,0.95,0.92}

\newcommand{\kw}[1]{{\color{magenta}\texttt{\bfseries\upshape #1}}}
\newcommand{\bluekw}[1]{{\color{blue}\texttt{\bfseries\upshape #1}}}

\lstdefinelanguage{Lex}{
  keywords={law, chapter, article, paragraph, point, rule, whenever, oblige, constitute, except, refine, observable, suppressable, hide, event, internal, predicate, type, is, enforceable, suppressing, causing, refinement, function, ONCE, NOT, SINCE, OR, AND, EXISTS},
  keywordstyle=\color{magenta}\bfseries,
  identifierstyle=\color{blue},
  sensitive=false,
  comment=[l]{\#},
  commentstyle=\color{purple}\ttfamily,
  stringstyle=\color{red}\ttfamily,
  morestring=[b]',
  morestring=[b]"
}

\lstdefinestyle{mystyle}{
  %backgroundcolor=\color{backcolour},   
  commentstyle=\color{codegreen},
  numberstyle=\tiny\color{white},
  stringstyle=\color{codepurple},
  basicstyle=\ttfamily\footnotesize,
  breakatwhitespace=false,         
  breaklines=true,                 
  captionpos=b,                    
  keepspaces=true,                 
  numbers=left,                    
  numbersep=5pt,                  
  showspaces=false,                
  showstringspaces=false,
  showtabs=false,                  
  tabsize=2
}

\lstset{language=Lex, style=mystyle}


% Begin document
\begin{document}

\begin{center}
    {\Large \textbf{Lex Cheat Sheet}}
\end{center}

\vspace{0.5em}

% Double column layout
\begin{multicols}{2}
  \small

\section{Labels}
\begin{tcolorbox}
  Available levels:

  \kw{law}, \kw{title}, \kw{chapter}, \kw{section}, \kw{article}, \kw{paragraph}, \kw{point}, \kw{subpoint}

  \bigskip

  Levels required for all rules:

  \kw{law}, \kw{article}
  

  \bigskip
  
  Example:
  \begin{lstlisting}
    law "Name of law"
    chapter "A" "Optional title of chapter"
    article "6" "Optional title of article"
    paragraph "1"
    point "1"
    ...
    article "7"
  \end{lstlisting}
  \vspace{-5mm}
\end{tcolorbox}

\section{Types}
\begin{tcolorbox}
  Base types:

  \kw{bool}, \kw{int}, \kw{float}, \kw{string}, \kw{span}, \kw{time}, \kw{money}\,\texttt{c}

  \bigskip

  Type declarations:
  \begin{lstlisting}
type user_consent_action
    
type user_id
    """a user identifier"""

type purpose is string
    """a purpose of processing"""
  \end{lstlisting}

  Possible values:
  \begin{lstlisting}
true, false                   # bool
1, 18, 42                     # int
1.2, 3.1415, 4.               # float
"hello", "Erika Mustermann"   # string
2s, 3m, 4h, 6d, 7M, 8y        # span
`2024-01-01 08:32:27`         # time
|CHF 2.|, |CHF 4000.25|       # money CHF
\end{lstlisting}
  \vspace{-3mm}
\end{tcolorbox}

\section{Events}
\begin{tcolorbox}
  Two keywords:

  \kw{event} (for actions), \kw{predicate} (for relations)

  \bigskip

  Capabilities:

  \kw{observable} $<$ \kw{suppressable}, \kw{causable}, \kw{causable observable},
  \kw{causable suppressable}, \kw{internal}

  \bigskip

  Non-\kw{internal} $\to$ `brute facts' in the world vs. \kw{internal} $\to$ `institutional facts'
  in the law.

  \bigskip

  Event declarations:
  \begin{lstlisting}
observable predicate PersonalData
    """data {d} is personal data of data subject {ds}"""
    d : data
    ds : data_subject

suppressable event ProcessesData
    """company {c} processes data {d} for purpose {p} as part of data processing activity {a}"""
    c : company
    d : data
    p : purpose
    a : activity

internal event IsLegal
    """processing activity {a} is legal"""
    a : activity
  \end{lstlisting}
    \vspace{-5mm}
\end{tcolorbox}

\section{Expressions}
\begin{tcolorbox}
  \begin{tabular}{ll}
    \texttt{( ... )} & Brackets \\
    \kw{true}, \kw{false}, \texttt{1}, \texttt{2.5}, ... & Constants \\
    \texttt{=}, \texttt{<>}, \texttt{<}, \texttt{>}, \texttt{<=}, \texttt{>=} & Relations\\
    \texttt{+}, \texttt{-}, \texttt{*}, \texttt{/}, \texttt{\^} & Arithmetics\\
    $\texttt{Event}\texttt{(t}_1\texttt{,}\dots\texttt{,t}_k\texttt{)}$ & Event/Predicate \\
    \kw{NOT} & Negation \\
    \kw{AND}, \kw{OR}, \kw{IMPLIES}, \kw{IFF} & Connectives \\
    \kw{EXISTS}~\texttt{x.} & There exists $x$...\\
    \kw{FORALL}~\texttt{x.} & For all $x$... \\
    \kw{ONCE} & `some time in the past...' \\
    \kw{HISTORICALLY} & `at all times in the past...'  \\
    \kw{EVENTUALLY} & `some time in the future...'\\
    \kw{ALWAYS} & `always in the future...'
  \end{tabular}

  \bigskip
  
  Temporal connectives can be followed by intervals \texttt{[a, b]}, \texttt{[a, b)},
  \texttt{(a, b]}, or \texttt{(a, b)} where \texttt{b} can be \texttt{*} ($\infty$).
  Bounds are integers followed by an optional time unit in \bluekw{s}, \bluekw{m},
  \bluekw{h}, \bluekw{d}, \bluekw{M}, \bluekw{Y} (default: \bluekw{s}).
\end{tcolorbox}

\section{Rules}
\begin{tcolorbox}
  Obligation rules: \emph{`whenever all of the premisses are true, then all the conclusions shall be true'}

  Syntax: \kw{whenever}... \kw{oblige}

  \begin{lstlisting}
# [GDPR Art. 5] Personal data shall be:
# 1. processed lawfully, fairly and in a
# transparent manner in relation to the
# data subject

article "5"
paragraph "1"
rule
    whenever
        ProcessesData(c, d, p, a)
    oblige
        IsLawful(a)
        IsFair(a)
        IsTransparent(a)
  \end{lstlisting}
  \vspace{-3mm}
  After this rule, we can add either of the two lines
  \begin{lstlisting}
enforceable suppressing ProcessesData
enforceable causing IsLawful, IsFair, IsTransparent    
  \end{lstlisting}
  \vspace{-5mm}
  \bigskip
  
  Constitutive rules: \emph{`whenever all the premisses are true, then all the conclusions are the case (by definition)'}

  The conclusions \textbf{must be} \kw{internal} events!

  Syntax: \kw{whenever}... \kw{constitute}

  \begin{lstlisting}
# Processing for a purpose is lawful when the 
# data subject has given consent for
# processing their data for that purpose
# (consent was given some time in the past)

rule
    whenever
        ProcessesData(c, d, p, a)
        PersonalData(d, ds)
        ONCE GivesConsent(ds, c, p)
    constitute
        IsLawful(a)
  \end{lstlisting}
  Exception rules: \emph{`whenever all the premises are true, then the rules in the conclusions shall not apply'}

  The conclusions \textbf{must be} labels!

  Syntax: \kw{whenever}... \kw{except}

  \begin{lstlisting}
# Art. 5 par. 1 shall not apply when there
# exists a legal warrant to use the data

rule
    whenever
        LegalWarrant(c, d)
    except
        article "5" paragraph "1"
  \end{lstlisting}
  \vspace{-5mm}
\end{tcolorbox}


\end{multicols}

\end{document}
%%% Local Variables:
%%% mode: latex
%%% TeX-master: t
%%% End:
